\begingroup
\tcbset{breakable/.default=false}
\chapter{质点运动的描述}

{
\graphicspath{{/Users/dongxingxi/Documents/物理竞赛一轮资料/sample/xes-original/kinematics-supplement/}{/Users/dongxingxi/Documents/物理竞赛一轮资料/sample/xes-original/kinematics-supplement/assets/}}
\section{理想物理模型：质点}
\newcommand{\EarthMoonScaleDiagram}{%
\begin{center}\begin{tikzpicture}[x=1cm,y=1cm,>=Latex]
  % 主图按同一线性比例绘制：中心距 10.5 cm，对应天体半径分别
  % 为 1.74 mm 和 0.48 mm。这样保留题图“真实尺度”的物理作用。
  \coordinate (E) at (-5.25,0);
  \coordinate (M) at (5.25,0);
  \shade[ball color=Blue!70] (E) circle (.174);
  \shade[ball color=TextGray!45] (M) circle (.048);
  \draw[<->,TextGray] (-5.25,-.45)--(5.25,-.45)
    node[midway,clearlabel,below=2pt]
      {平均中心距 $3.84\times10^5\,\mathrm{km}$};
  \draw[TextGray] (E)--++(-.42,.58);
  \node[clearlabel,above right,align=left] at (-5.72,.55)
    {地球 $R_E\approx6.37\times10^3\,\mathrm{km}$\\$M_E\approx5.97\times10^{24}\,\mathrm{kg}$};
  \draw[TextGray] (M)--++(.35,.58);
  \node[clearlabel,above left,align=right] at (5.67,.55)
    {月球 $R_M\approx1.74\times10^3\,\mathrm{km}$\\$M_M\approx7.35\times10^{22}\,\mathrm{kg}$};
  \node[clearlabel,below] at (0,-1.05)
    {半径与中心距采用同一比例；$M_M/M_E\approx0.0123$};
\end{tikzpicture}\figcap{月球、地球及月地距离的真实比例。研究月球绕地球公转时可把月球近似为质点；研究着陆时则不能。}\label{ch02:fig:k01-2-1}\end{center}}


自然界中的物体都处在运动过程中；所谓静止，只能相对于某个指定参考系而言。在经典力学的运动学描述中，不把“绝对静止”作为可直接判定的状态。学习力学的第一步，是建立一套能够准确描写运动的方法：先说明研究对象，再规定参考系和坐标，最后用位置随时间的变化表达运动。只有在运动学描述清楚之后，才有可能进一步研究运动状态变化的原因。

物理研究通常从最简单的对象开始。把一个具有确定质量、但形状和内部结构可以忽略的运动物体抽象成几何点，称为\textbf{质点}（mass point）。质点不是客观存在的“无大小物体”，而是针对具体问题建立的理想模型。

能否把物体看成质点，不能只比较物体自身大小，还要比较它的大小与研究尺度，并判断姿态、转动和形变是否影响所问物理量。例如：

\begin{itemize}
  \item 研究跳水运动员入水姿态时，身体各部分的位置关系直接影响结果，不能把运动员看成质点；研究他从北京到上海的行程时，身体尺寸相对于路程可以忽略，可用质点表示。
  \item 研究月球绕地球公转时，月球半径远小于地月距离，通常可把月球当成质点；研究探测器在月面着陆时，月球表面的曲率、地形和自转均可能重要，不能再用同一个简化。
  \item 研究乒乓球质心的抛体轨迹时可以暂时忽略球的大小；研究旋转引起的弧线或球与球拍接触时，必须保留半径和自转。
\end{itemize}

\EarthMoonScaleDiagram

建模时应明确写出“忽略了什么”。质点模型通常忽略物体的形状、姿态、转动和内部自由度，只保留质量与位置。若后续推导需要角速度、转动惯量、接触几何或形变能，就说明质点模型已经不足，必须升级为刚体、连续介质或可变形体模型。

判断模型是否合理的可靠方法，是做尺度和极限检查。设物体特征尺寸为 $\ell$、运动轨迹的特征尺度为 $L$。若 $\ell/L\ll1$，并且观测精度不要求分辨 $\ell$ 量级的差异，质点近似通常可靠；若结果对 $\ell$ 的一阶修正已经不可忽略，就不能仅凭“看起来很小”继续使用质点模型。
}

{
\graphicspath{{/Users/dongxingxi/Documents/物理竞赛一轮资料/sample/xes-original/kinematics-supplement/}{/Users/dongxingxi/Documents/物理竞赛一轮资料/sample/xes-original/kinematics-supplement/assets/}}
\section{时间、位移、速度与加速度}
\newcommand{\SprintGraph}{%
\begin{center}\begin{tikzpicture}
\begin{axis}[
  width=.78\linewidth,height=6.0cm,
  axis lines=left,xlabel={$t/\mathrm{s}$},ylabel={$x/\mathrm{m}$},
  xmin=0,xmax=12.5,ymin=0,ymax=105,
  xtick={0,...,12},ytick={0,20,...,100},
  grid=major,grid style={draw=RuleGray},
  every axis label/.style={fill=none},tick label style={font=\small}
]
\addplot[kcurve,mark=*,mark size=1.3pt] coordinates {
  (0,0)(1,1)(2,4)(3,9)(4,16)(5,25)(6,35)(7,45)(8,55)(9,65)(10,75)(11,85)(12,100)
};
\end{axis}
\end{tikzpicture}\figcap{短跑选手位置随时间变化。曲线越陡，位置变化越快。}\label{ch02:fig:k02-2-2}\end{center}}

\newcommand{\PositionSlopeDiagram}{%
\begin{center}\begin{tikzpicture}[scale=.94,>=Latex]
  \foreach \dx/\dy/\mode in {-3/0/1,3/0/2,0/-5.2/3}{
    \begin{scope}[shift={(\dx,\dy)}]
      \draw[kaxis] (-1.8,0)--(1.9,0) node[right,clearlabel]{$t$};
      \draw[kaxis] (-1.55,-1.35)--(-1.55,1.55) node[above,clearlabel]{$x$};
      \ifnum\mode=1
        \draw[kcurve] (-1.35,-.9)--(1.55,1.05);
        \node[clearlabel,below,align=center] at (0,-1.55){匀速运动\\直线斜率恒定};
      \else\ifnum\mode=2
        \draw[kcurve] plot[smooth] coordinates{(-1.35,.12)(-.7,.35)(0,.64)(.7,1.0)(1.55,1.27)};
        \coordinate (A) at (-.7,.35);\coordinate (B) at (.7,1.0);
        \fill (A) circle (.04);
        \fill (B) circle (.04);
        \draw[secant] (A)--(B);
        \draw[kguide] (-.7,0)--(A) (.7,0)--(B);
        \draw[kguide] (-1.55,.35)--(A) (-1.55,1.0)--(B);
        \draw[<->,TextGray] (-.7,-.45)--(.7,-.45)
          node[midway,clearlabel,below=1pt]{$\Delta t$};
        \draw[<->,TextGray] (1.18,.35)--(1.18,1.0)
          node[midway,clearlabel,right=1pt]{$\Delta x$};
        \draw[TextGray] (A)--(.38,.35);
        \draw[TextGray] (-.30,.35) arc[start angle=0,end angle=24.90,radius=.40];
        \node[clearlabel] at (.20,.58){$\alpha$};
        \node[clearlabel,above left=2pt] at (A) {$A$};
        \node[clearlabel,above right=2pt] at (B) {$B$};
        \node[clearlabel,below=1pt] at (-.7,0) {$t_1$};
        \node[clearlabel,below=1pt] at (.7,0) {$t_2$};
        \node[clearlabel,left=1pt] at (-1.55,.35) {$x$};
        \node[clearlabel,left=1pt] at (-1.55,1.0) {$x+\Delta x$};
        \node[clearlabel,below,align=center] at (0,-1.55){割线斜率：平均速度\\$\Delta t=t_2-t_1$};
      \else
        \draw[kcurve,domain=-1.35:1.5,samples=80] plot (\x,{.12+.65*\x-.1*\x*\x});
        \coordinate (A) at (0,.12);\fill (A) circle (.04);
        \draw[tangent] (-1.0,-.53)--(1.1,.835);
        \node[clearlabel,below right=3pt] at (A) {$A$};
        \node[clearlabel,below,align=center] at (0,-1.55){切线斜率\\瞬时速度};
      \fi\fi
    \end{scope}
  }
\end{tikzpicture}\figcap{$x$-$t$ 图像中平均速度与瞬时速度的几何意义。}\label{ch02:fig:k02-2-3}\end{center}}

\newcommand{\VelocityAreaDiagram}{%
\begin{center}\begin{tikzpicture}[scale=.82,>=Latex]
  \foreach \dx/\mode in {-4.7/1,0/2,4.7/3}{
    \begin{scope}[xshift=\dx cm]
      \draw[kaxis] (-1.8,0)--(1.9,0) node[right,clearlabel]{$t$};
      \draw[kaxis] (-1.55,-.15)--(-1.55,1.65) node[above,clearlabel]{$v$};
      \ifnum\mode=1
        \fill[Blue!12] (-1.1,0) rectangle (1.3,.9);
        \draw[kcurve] (-1.25,.9)--(1.5,.9);
        \draw[kguide] (-1.1,0)--(-1.1,.9) (1.3,0)--(1.3,.9);
        \node[clearlabel,below=2pt] at (-1.1,0){$A$};
        \node[clearlabel,below=2pt] at (1.3,0){$B$};
        \node[clearlabel,above left=1pt] at (-1.1,.9){$C$};
        \node[clearlabel,above right=1pt] at (1.3,.9){$D$};
        \node[clearlabel,below] at (0,-.55){匀速：矩形面积};
      \else
        \path[fill=Blue!12] (-1.1,0)--
          plot[domain=-1.1:1.3,samples=60]
            ({\x},{0.55+0.75*(\x+1.1)-0.22*(\x+1.1)^2})--(1.3,0)--cycle;
        \draw[kcurve]
          plot[domain=-1.1:1.3,samples=60]
            ({\x},{0.55+0.75*(\x+1.1)-0.22*(\x+1.1)^2});
        \draw[kguide] (-1.1,0)--(-1.1,.55) (1.3,0)--(1.3,1.08);
        \node[clearlabel,below=2pt] at (-1.1,0){$A$};
        \node[clearlabel,below=2pt] at (1.3,0){$B$};
        \node[clearlabel,above left=1pt] at (-1.1,.55){$C$};
        \node[clearlabel,above right=1pt] at (1.3,1.08){$D$};
        \ifnum\mode=2
          \foreach \x in {-1.05,-.85,...,1.25}{
            \draw[Blue!65] (\x,0)--
              (\x,{0.55+0.75*(\x+1.1)-0.22*(\x+1.1)^2});}
          \node[clearlabel,below] at (0,-.55){窄矩形求和};
        \else
          \node[clearlabel,below] at (0,-.55){极限：曲线下面积};
        \fi
      \fi
    \end{scope}
  }
\end{tikzpicture}\figcap{$v$-$t$ 图像下的有向面积等于位移；若求路程，应对 $|v|$ 积分。}\label{ch02:fig:k02-2-4}\end{center}}

\newcommand{\ExampleTwoSixGraphs}{%
\begin{center}\begin{tikzpicture}[scale=.78,>=Latex]
  \foreach \dx/\mode in {-4.6/1,0/2,4.6/3}{
    \begin{scope}[xshift=\dx cm]
      \draw[kaxis] (-1.65,0)--(1.75,0) node[right,clearlabel]{$t$};
      \draw[kaxis] (-1.45,-.55)--(-1.45,2.05) node[above,clearlabel]{$x$};
      \ifnum\mode=1
        \draw[kcurve] (-1.35,0)--(0,1.35)--(1.35,0);
        \draw[vector] (-.92,.43)--(-.42,.93);
        \draw[vector] (.42,.93)--(.92,.43);
        \node[clearlabel,below] at (0,-.62){先正向、后反向};
      \else\ifnum\mode=2
        \draw[kcurve,domain=-1.35:1.30,samples=80] plot (\x,{.25*(\x+1.35)^2});
        \draw[vector] (-.65,.12)--(-.15,.34);
        \draw[vector] (.35,.73)--(.82,1.25);
        \node[clearlabel,below] at (0,-.62){斜率为正且增大};
      \else
        \draw[kcurve,domain=-1.35:1.35,samples=100] plot (\x,{1.45*sqrt(max(0,1-(\x/1.35)^2))});
        \draw[vector] (-.85,.99)--(-.35,1.34);
        \draw[vector] (.35,1.34)--(.85,.99);
        \node[clearlabel,below] at (0,-.62){斜率由正变为负};
      \fi\fi
    \end{scope}}
\end{tikzpicture}\figcap{三幅 $x$-$t$ 图的切线斜率分别给出速度的方向与变化趋势。}\label{ch02:fig:k02-ex-2-6}\end{center}}

\newcommand{\ExampleTwoSevenGraphs}{%
\begin{center}\begin{tikzpicture}[scale=.9,>=Latex]
  \begin{scope}[shift={(-4.9,-1.1)}]
    \draw[kaxis] (0,0)--(2.75,0) node[right,clearlabel]{$t$};
    \draw[kaxis] (0,-1.55)--(0,1.85) node[above,clearlabel]{$v$};
    \draw[kcurve] (0,1.3)--(2.4,-1.3);
    \node[clearlabel,left] at (0,1.3){$v_0$};
    \node[clearlabel,left] at (0,-1.3){$-v_0$};
    \draw[kguide] (0,-1.3)--(2.4,-1.3);
    \node[clearlabel,below left] at (0,0){$O$};
    \node[clearlabel,above right=2pt] at (1.2,0){$T/2$};
    \node[clearlabel,above] at (2.4,0){$T$};
    \node[clearlabel,below] at (1.2,-1.7){已知速度图};
  \end{scope}
  \foreach \dx/\dy/\mode/\tag in {-.7/0/1/A,3.3/0/2/B,-.7/-3.3/3/C,3.3/-3.3/4/D}{
    \begin{scope}[shift={(\dx,\dy)}]
      \draw[kaxis] (0,0)--(2.7,0) node[right,clearlabel]{$t$};
      \draw[kaxis] (0,-.25)--(0,1.65) node[above,clearlabel]{$x$};
      \ifnum\mode=1
        \draw[kcurve] (0,0)--(1.2,1.15)--(2.4,0);
      \else\ifnum\mode=2
        \draw[kcurve,domain=0:1,samples=70]
          plot ({2.4*\x},{4.6*\x*(1-\x)});
      \else\ifnum\mode=3
        \draw[kcurve,domain=0:1,samples=70] plot ({2.4*\x},{1.15*(2*\x-\x*\x)});
      \else
        \draw[kcurve] (0,1.15)--(1.2,0)--(2.4,.85);
      \fi\fi\fi
      \node[clearlabel,below left] at (0,0){$O$};
      \node[clearlabel,below=3pt] at (1.2,0){$T/2$};
      \node[clearlabel,below=3pt] at (2.4,0){$T$};
      \node[clearlabel,below] at (1.2,-.65){$\tag$};
    \end{scope}}
\end{tikzpicture}\figcap{左图为给定的 $v$-$t$ 图，右侧四图为待判断的位移图；正确图必须连续且斜率等于给定速度。}\label{ch02:fig:k02-ex-2-7}\end{center}}

\newcommand{\ExampleTwoEightGraphs}{%
\begin{center}\begin{tikzpicture}[scale=.78,>=Latex]
  \foreach \dx/\mode in {-4.4/1,0/2,4.4/3}{
    \begin{scope}[xshift=\dx cm]
      \draw[kaxis] (-1.55,0)--(1.75,0) node[right,clearlabel]{时间};
      \draw[kaxis] (-1.35,-.25)--(-1.35,1.85) node[above,clearlabel]{$v$};
      \ifnum\mode=1
        \fill[Blue!12] (-1.25,0)--(1.25,0)--(1.25,1.35)--cycle;
        \draw[kcurve] (-1.25,0)--(1.25,1.35);
        \draw[kguide] (1.25,0)--(1.25,1.35);
        \draw[kguide] (-1.35,1.35)--(1.25,1.35);
        \node[clearlabel,left] at (-1.35,1.35){$v$};
      \else\ifnum\mode=2
        \fill[Blue!12] (-1.25,0)--(-1.25,.55)--(1.25,1.35)--(1.25,0)--cycle;
        \draw[kcurve] (-1.25,.55)--(1.25,1.35);
        \draw[kguide] (1.25,0)--(1.25,1.35);
        \draw[kguide] (-1.35,.55)--(-1.25,.55) (-1.35,1.35)--(1.25,1.35);
        \node[clearlabel,left] at (-1.35,.55){$v_0$};
        \node[clearlabel,left] at (-1.35,1.35){$v_1$};
      \else
        \fill[Blue!12] (-1.25,0)--(-.42,1.35)--(1.25,0)--cycle;
        \draw[kcurve] (-1.25,0)--(-.42,1.35)--(1.25,0);
        \draw[kguide] (-.42,0)--(-.42,1.35);
        \draw[kguide] (-1.35,1.35)--(-.42,1.35);
        \node[clearlabel,left] at (-1.35,1.35){$v_0$};
        \node[clearlabel,below] at (-.42,0){$t/3$};
      \fi\fi
      \node[clearlabel,below] at (1.25,0){$t$};
    \end{scope}}
\end{tikzpicture}\figcap{总路程等于三幅 $v$-$t$ 图中阴影面积；梯形或三角形面积直接给出结果。}\label{ch02:fig:k02-ex-2-8}\end{center}}

\newcommand{\FreeFallMeetingDiagram}{%
\begin{center}\begin{tikzpicture}[scale=.9,>=Latex]
  \draw[very thick] (-1.4,0)--(1.4,0);
  \node[point,label={[clearlabel]below right=3pt:$B$}] (B) at (0,0) {};
  \node[point,label={[clearlabel]left:$A$}] (A) at (0,3.2) {};
  \draw[<->,TextGray] (.8,0)--(.8,3.2) node[midway,clearlabel,right]{$h$};
  \draw[DeepBlue,very thick,->] (B)--(0,1.15) node[midway,clearlabel,left]{$v_0$};
  \draw[Blue,very thick,->] (A)--(0,2.45) node[midway,clearlabel,left]{$g$};
\end{tikzpicture}\figcap{解图：两球具有相同的重力加速度，因此相对加速度为零，初始间距以相对速度 $v_0$ 关闭。}\label{ch02:fig:k02-ex-2-17}\end{center}}

\newcommand{\BounceDiagram}{%
\begin{center}\begin{tikzpicture}[scale=.85,>=Latex]
  \draw[very thick] (-1.7,0)--(1.7,0);
  \fill (0,2.7) circle (.12);
  \fill (0,.12) circle (.12);
  \draw[kguide] (0,.27)--(0,2.55);
  \draw[<->,TextGray] (.88,0)--(.88,2.7)
    node[midway,clearlabel,right=2pt]{$H$};
  \draw[DeepBlue,very thick,->] (-.35,1.0)--(-.35,.25) node[midway,clearlabel,left]{$u$};
  \draw[Blue,very thick,->] (.35,.25)--(.35,1.0) node[midway,clearlabel,right]{$eu$};
\end{tikzpicture}\figcap{解图：每次反弹速度按 $e$ 成等比衰减，飞行时间按 $e$ 衰减，路程按 $e^2$ 衰减。}\label{ch02:fig:k02-ex-2-19}\end{center}}

\newcommand{\KQuestionAB}{%
\begin{center}\begin{minipage}{.97\linewidth}\centering
\begin{tikzpicture}[>=Latex,scale=.9]

\draw[kaxis] (-.25,0)--(6.6,0) node[right]{$t/\mathrm s$};
\draw[kaxis] (0,-2.7)--(0,1.2) node[above]{$x/\mathrm m$};
\foreach \t in {1,...,6}{\draw (\t,0)--(\t,.08);\node[above=2pt] at (\t,0){\t};}
\foreach \y/\v in {-2.4/-60,-1.6/-40,-.8/-20,.8/20}{\draw (0,\y)--(.08,\y);\node[left] at (0,\y){$\v$};}
\node[below left] at (0,0){$O$};

\end{tikzpicture}
\figcap{题图：按表中数据在空坐标系中作图。}\label{ch02:ref:question-figure-2-1}
\end{minipage}\end{center}
}

\newcommand{\KQuestionAC}{%
\begin{center}\begin{minipage}{.97\linewidth}\centering
\begin{tikzpicture}[>=Latex,scale=.85]

\foreach \dx/\i in {0/1,4.3/2,8.6/3}{\begin{scope}[xshift=\dx cm]
\draw[kaxis] (-.2,0)--(3.1,0) node[right]{$t$};
\draw[kaxis] (0,-.6)--(0,1.8) node[above]{$x$};
\node[below left] at (0,0){$O$};\node at (1.4,-.85){(\i)};
\end{scope}}

\end{tikzpicture}
\figcap{题图：分别画出三个运动函数的图像。}\label{ch02:ref:question-figure-2-2}
\end{minipage}\end{center}
}

\newcommand{\KQuestionAD}{%
\begin{center}\begin{minipage}{.97\linewidth}\centering
\begin{tikzpicture}[>=Latex,scale=.9]

\foreach \dx/\mode in {0/1,6.2/2}{\begin{scope}[xshift=\dx cm]
\draw[kaxis] (-.2,0)--(4.6,0) node[right]{$t$};
\draw[kaxis] (0,-1.25)--(0,1.3) node[above]{$x$};
\ifnum\mode=1
\draw[kcurve,domain=0:4.25,samples=200] plot(\x,{.8*sin(280*\x)});
\else
\draw[kcurve,domain=0:4.25,samples=240] plot(\x,{1.05*exp(-.85*\x)*cos(450*\x)});
\fi
\node[below left] at (0,0){$O$};\node at (2,-1.6){(\mode)};
\end{scope}}

\end{tikzpicture}
\figcap{题图：两个质点的一维位置随时间变化，坐标轴未给定数值比例。}\label{ch02:ref:question-figure-2-3}
\end{minipage}\end{center}
}

\newcommand{\KQuestionAG}{%
%
\begin{center}\begin{tikzpicture}[scale=.78,>=Latex]
  \foreach \dx/\mode in {-4.6/1,0/2,4.6/3}{
    \begin{scope}[xshift=\dx cm]
      \draw[kaxis] (-1.65,0)--(1.75,0) node[right,clearlabel]{$t$};
      \draw[kaxis] (-1.45,-.55)--(-1.45,2.05) node[above,clearlabel]{$x$};
      \ifnum\mode=1
        \draw[kcurve] (-1.35,0)--(0,1.35)--(1.35,0);
      \else\ifnum\mode=2
        \draw[kcurve,domain=-1.35:1.30,samples=80] plot (\x,{.25*(\x+1.35)^2});
      \else
        \draw[kcurve,domain=-1.35:1.35,samples=100] plot (\x,{1.45*sqrt(max(0,1-(\x/1.35)^2))});
      \fi\fi
    \end{scope}}
\end{tikzpicture}\figcap{题图：从左到右为三幅给定的位置—时间图像。}\label{ch02:ref:question-figure-2-6}\end{center}
}

\newcommand{\KQuestionAH}{%
%
\begin{center}\begin{tikzpicture}[scale=.9,>=Latex]
  \begin{scope}[shift={(-4.9,-1.1)}]
    \draw[kaxis] (0,0)--(2.75,0) node[right,clearlabel]{$t$};
    \draw[kaxis] (0,-1.55)--(0,1.85) node[above,clearlabel]{$v$};
    \draw[kcurve] (0,1.3)--(2.4,-1.3);
    \node[clearlabel,left] at (0,1.3){$v_0$};
    \node[clearlabel,left] at (0,-1.3){$-v_0$};
    \draw[kguide] (0,-1.3)--(2.4,-1.3);
    \node[clearlabel,below left] at (0,0){$O$};
    \node[clearlabel,above right=2pt] at (1.2,0){$T/2$};
    \node[clearlabel,above] at (2.4,0){$T$};
    \node[clearlabel,below] at (1.2,-1.7){已知速度图};
  \end{scope}
  \foreach \dx/\dy/\mode/\tag in {-.7/0/1/A,3.3/0/2/B,-.7/-3.3/3/C,3.3/-3.3/4/D}{
    \begin{scope}[shift={(\dx,\dy)}]
      \draw[kaxis] (0,0)--(2.7,0) node[right,clearlabel]{$t$};
      \draw[kaxis] (0,-.25)--(0,1.65) node[above,clearlabel]{$x$};
      \ifnum\mode=1
        \draw[kcurve] (0,0)--(1.2,1.15)--(2.4,0);
      \else\ifnum\mode=2
        \draw[kcurve,domain=0:1,samples=70]
          plot ({2.4*\x},{4.6*\x*(1-\x)});
      \else\ifnum\mode=3
        \draw[kcurve,domain=0:1,samples=70] plot ({2.4*\x},{1.15*(2*\x-\x*\x)});
      \else
        \draw[kcurve] (0,1.15)--(1.2,0)--(2.4,.85);
      \fi\fi\fi
      \node[clearlabel,below left] at (0,0){$O$};
      \node[clearlabel,below=3pt] at (1.2,0){$T/2$};
      \node[clearlabel,below=3pt] at (2.4,0){$T$};
      \node[clearlabel,below] at (1.2,-.65){$\tag$};
    \end{scope}}
\end{tikzpicture}\figcap{题图：左图给出速度随时间的变化，右侧 A—D 为待选的位置图像。}\label{ch02:ref:question-figure-2-7}\end{center}
}

\newcommand{\KQuestionAI}{%
%
\begin{center}\begin{tikzpicture}[scale=.78,>=Latex]
  \foreach \dx/\mode in {-4.4/1,0/2,4.4/3}{
    \begin{scope}[xshift=\dx cm]
      \draw[kaxis] (-1.55,0)--(1.75,0) node[right,clearlabel]{时间};
      \draw[kaxis] (-1.35,-.25)--(-1.35,1.85) node[above,clearlabel]{$v$};
      \ifnum\mode=1
        \draw[kcurve] (-1.25,0)--(1.25,1.35);
        \draw[kguide] (1.25,0)--(1.25,1.35);
        \draw[kguide] (-1.35,1.35)--(1.25,1.35);
        \node[clearlabel,left] at (-1.35,1.35){$v$};
      \else\ifnum\mode=2
        \draw[kcurve] (-1.25,.55)--(1.25,1.35);
        \draw[kguide] (1.25,0)--(1.25,1.35);
        \draw[kguide] (-1.35,.55)--(-1.25,.55) (-1.35,1.35)--(1.25,1.35);
        \node[clearlabel,left] at (-1.35,.55){$v_0$};
        \node[clearlabel,left] at (-1.35,1.35){$v_1$};
      \else
        \draw[kcurve] (-1.25,0)--(-.42,1.35)--(1.25,0);
        \draw[kguide] (-.42,0)--(-.42,1.35);
        \draw[kguide] (-1.35,1.35)--(-.42,1.35);
        \node[clearlabel,left] at (-1.35,1.35){$v_0$};
        \node[clearlabel,below] at (-.42,0){$t/3$};
      \fi\fi
      \node[clearlabel,below] at (1.25,0){$t$};
    \end{scope}}
\end{tikzpicture}\figcap{题图：三个运动过程的速度—时间图像，求相应时间内的路程。}\label{ch02:ref:question-figure-2-8}\end{center}
}

\newcommand{\KQuestionAR}{%
%
\begin{center}\begin{tikzpicture}[scale=.9,>=Latex]
  \draw[very thick] (-1.4,0)--(1.4,0);
  \node[point,label={[clearlabel]below right=3pt:$B$}] (B) at (0,0) {};
  \node[point,label={[clearlabel]left:$A$}] (A) at (0,3.2) {};
  \draw[<->,TextGray] (.8,0)--(.8,3.2) node[midway,clearlabel,right]{$h$};
  \draw[DeepBlue,very thick,->] (B)--(0,1.15) node[midway,clearlabel,left]{$v_0$};
  \draw[Blue,very thick,->] (A)--(0,2.45) node[midway,clearlabel,left]{$g$};
\end{tikzpicture}\figcap{题图：A 球从上方释放，B 球从下方以初速度 $v_0$ 竖直抛出，初始间距为 $h$。}\label{ch02:ref:question-figure-2-17}\end{center}
}

\newcommand{\KQuestionAT}{%
%
\begin{center}\begin{tikzpicture}[scale=.85,>=Latex]
  \draw[very thick] (-1.7,0)--(1.7,0);
  \fill (0,2.7) circle (.12);
  \fill (0,.12) circle (.12);
  \draw[kguide] (0,.27)--(0,2.55);
  \draw[<->,TextGray] (.88,0)--(.88,2.7)
    node[midway,clearlabel,right=2pt]{$H$};
  \draw[DeepBlue,very thick,->] (-.35,1.0)--(-.35,.25) node[midway,clearlabel,left]{$u$};
  \draw[Blue,very thick,->] (.35,.25)--(.35,1.0) node[midway,clearlabel,right]{$eu$};
\end{tikzpicture}\figcap{题图：初始高度为 $H$；一次碰撞前、后的速率分别用 $u$、$eu$ 表示。}\label{ch02:ref:question-figure-2-19}\end{center}
}


把运动物体抽象为质点以后，质点的位置代表整个物体的位置。我们从一维直线运动开始：在运动直线上取原点、正方向和长度单位，建立数轴 $Ox$。质点在时刻 $t$ 的位置由坐标 $x(t)$ 唯一确定。坐标的符号表示点位于原点的哪一侧，并不表示“路程为负”。

这种简化只是描写运动的第一步。即使研究对象已经缩成一个点，要完整说明它怎样运动，仍须知道它在每一时刻处在什么位置。先从直线运动建立方法，再推广到曲线运动和真实物体的复杂运动，是本章的基本顺序；等运动本身被描述清楚之后，才进一步追问支配运动变化的原因。

建立一维坐标系时，原点、正方向和长度单位都可按问题需要选择，通常以米或千米为长度单位。质点在某时刻的坐标又称它在该坐标系中的位置。改变原点或正方向会改变坐标的数值或符号，却不会改变实际运动。若已知 $t=0,1,2,\ldots$ 各时刻的坐标，就已掌握一部分运动信息；记录时刻越密，所知道的运动细节越多，但只有连续的函数才能原则上给出每一时刻的位置。

描写同一运动可以采用叙述、列表、图像和函数四种方式。叙述便于日常交流；列表适合记录实验数据；图像最直观；函数则能把连续运动在任意时刻的位置完整表达出来。以下短跑数据同时说明四种表达之间的联系。

若用叙述表达表~\ref{ch02:tab:k02-2-1}中的比赛过程，就要依次说出“第 $1\,\mathrm{s}$ 跑到 $1\,\mathrm{m}$、第 $2\,\mathrm{s}$ 跑到 $4\,\mathrm{m}$、……、第 $12\,\mathrm{s}$ 到达 $100\,\mathrm{m}$”。这种说法能够传递信息，却繁琐且不容易立即比较各阶段的快慢。列表节省空间，适合实验记录，但仍不能一眼看出变化趋势。把时间作为横轴、离开起点的距离作为纵轴，数据点连成图~\ref{ch02:fig:k02-2-2}，比赛开始阶段曲线越来越陡，说明选手在加速；中段斜率近似不变；终点前曲线又略变陡，可看出冲刺过程。

\begin{center}\small
\begin{tabular}{c*{12}{r}}
\toprule
$t/\mathrm{s}$&1&2&3&4&5&6&7&8&9&10&11&12\\
$x/\mathrm{m}$&1&4&9&16&25&35&45&55&65&75&85&100\\
\bottomrule
\end{tabular}
\booktablecaption{不同时刻短跑选手的位置。}\label{ch02:tab:k02-2-1}
\end{center}

\SprintGraph

若函数 $x=x(t)$ 已知，就能在定义域内求任意时刻的位置。图像并不是质点在空间中的轨迹，而是“位置”这个物理量随“时间”变化的关系图；横轴和纵轴的物理意义不能混淆。

从实验得到函数通常要经历“测量—作图—猜测规律—检验”的过程。有限个数据点本身并不能保证点间一定按某条光滑曲线运动；把点连起来，实际上加入了质点连续运动以及测量误差较小等物理判断。若在某段出现突然跳点，应先检查是否发生碰撞、换向或记录错误，而不能为了图线好看强行作光滑连接。反过来，给定解析函数后可以任意加密取点，所画曲线只是函数关系的可视化。

还要区分“时刻”和“时间间隔”。$t$ 是相对所选计时零点的读数，改变计时起点会改变它的数值；$\Delta t=t_2-t_1$ 是两个事件间隔，通常不随计时零点改变。位置 $x$ 依赖坐标原点，位置变化 $\Delta x=x_2-x_1$ 则与平移坐标原点无关。许多运动公式实际联系的是间隔和变化量，列式前把初始时刻平移到 $t=0$ 只是方便，并非物理规律只在零时刻开始成立。

叙述、表格和有限个实验点都不能完全表现运动的连续性。若在两个已记录时刻之间质点仍连续运动，就还存在无穷多个没有列出的时刻。函数表达式
\begin{equation}
  x=x(t)
\end{equation}
把任意允许时刻映射为一个位置，因此最适合精确分析。例如 $x=3t$ 表示相等时间内坐标增加量相等；$x=5t^2$ 的图线随时间变陡，表示相同时间内的坐标增加量逐渐增大。函数和图像各有所长：函数便于定量运算，图像便于观察单调性、转折、周期和极值。

\begin{origexamplebox}{由测量数据作位置图像}\label{ch02:ex:k02-2-1}
建立竖直向上的坐标系表示从原点竖直上抛物体的位置。实验在 $t=0,0.5,1,\ldots,6\,\mathrm{s}$ 测得位置依次为
\[
0,\ 8.75,\ 15,\ 18.75,\ 20,\ 18.75,\ 15,\ 8.75,\ 0,\ -11.25,\ -25,\ -41.25,\ -60\ \mathrm{m}.
\]
在坐标系中标出数据并连接成光滑曲线。
\par 题图见图~\ref{ch02:ref:question-figure-2-1}。
\KQuestionAB
\end{origexamplebox}
\Needspace{6\baselineskip}\begin{xesolution}
\textbf{第一步：确定坐标轴和比例。}
横轴表示时刻 $t$，范围取 $0$ 到 $6\,\mathrm{s}$；纵轴表示相对所选原点的竖直坐标 $x$，
范围至少要覆盖 $-60\,\mathrm{m}$ 到 $20\,\mathrm{m}$。把题中每一对 $(t,x)$ 作为一个点标在
$x$--$t$ 坐标系中。这里画的是“位置随时间变化”的图像，不是物体在竖直空间中走出的轨迹。

\textbf{第二步：由数据判断最简单的函数形式。}
相邻测量时刻的间隔均为
\[
 h=0.5\,\mathrm{s}.
\]
相邻位置的一阶差分依次为
\[
\begin{aligned}
\Delta x_i={}&8.75,\ 6.25,\ 3.75,\ 1.25,\ -1.25,\ -3.75,\\
              &-6.25,\ -8.75,\ -11.25,\ -13.75,\ -16.25,\ -18.75
              \quad(\mathrm m).
\end{aligned}
\]
这些数每次都减少 $2.50\,\mathrm m$，所以二阶差分为常量
\[
 \Delta^2x_i=x_{i+2}-2x_{i+1}+x_i=-2.50\,\mathrm m.
\]
等时间间隔数据具有恒定二阶差分，最简单的光滑模型是二次函数
\[
 x(t)=At^2+Bt+C.
\]
对二次函数，步长为 $h$ 时有
\[
 x(t+2h)-2x(t+h)+x(t)=2Ah^2.
\]
代入 $h=0.5\,\mathrm s$ 和二阶差分 $-2.50\,\mathrm m$，得到
\[
 -2.50\,\mathrm m=2A(0.5\,\mathrm s)^2,
 \qquad A=-5.0\,\mathrm{m/s^2}.
\]
又由 $x(0)=0$ 得 $C=0$；把 $x(0.5\,\mathrm s)=8.75\,\mathrm m$ 代入，
\[
 8.75\,\mathrm m
 =(-5.0\,\mathrm{m/s^2})(0.5\,\mathrm s)^2
 +B(0.5\,\mathrm s),
\]
解得 $B=20\,\mathrm{m/s}$。因此所有测量点都满足
\[
 \boxed{x(t)=(20\,\mathrm{m/s})t-(5.0\,\mathrm{m/s^2})t^2}.
\]

\textbf{第三步：确定图像的关键点和运动含义。}
把二次函数配方：
\[
 x(t)=20\,\mathrm m
 -(5.0\,\mathrm{m/s^2})(t-2.0\,\mathrm s)^2.
\]
所以图像是开口向下的抛物线，对称轴为 $t=2.0\,\mathrm s$，顶点为
$(2.0\,\mathrm s,20\,\mathrm m)$。令 $x=0$，得到 $t=0$ 和 $t=4.0\,\mathrm s$：物体
在 $4.0\,\mathrm s$ 时回到坐标原点。$t>4.0\,\mathrm s$ 后 $x<0$，表示物体位于所选原点下方，
并不是“路程为负”。

图线斜率给出速度。由
\[
 v(t)=\frac{\dd x}{\dd t}
 =20\,\mathrm{m/s}-(10\,\mathrm{m/s^2})t
\]
可见：$0<t<2\,\mathrm s$ 时斜率为正，物体上升；$t=2\,\mathrm s$ 时斜率为零，到达最高点；
$t>2\,\mathrm s$ 时斜率为负，物体下降。加速度恒为
$a=-10\,\mathrm{m/s^2}$，与竖直向上为正的坐标约定一致。量纲检查中，两项
$(20\,\mathrm{m/s})t$ 与 $(5.0\,\mathrm{m/s^2})t^2$ 都是长度。

最后需要说明：有限个测量点本身不能唯一确定点与点之间的真实运动。这里所有点具有严格恒定的二阶差分，
因此采用匀加速度对应的二次函数作最简单的光滑连接；若没有这一物理建模假设，还可以画出无穷多条经过同一组数据点的光滑曲线。
\begin{center}\begin{tikzpicture}
\begin{axis}[width=.72\linewidth,height=5.8cm,axis lines=middle,
  xlabel={$t/\mathrm{s}$},ylabel={$x/\mathrm{m}$},
  xlabel style={at={(axis description cs:1.01,.49)},anchor=west},
  ylabel style={rotate=0,at={(axis description cs:.035,1.02)},anchor=south west},
  xmin=0,xmax=6.2,ymin=-65,ymax=25,grid=major,grid style={draw=RuleGray},
  every axis label/.style={fill=none}]
\addplot[kcurve,domain=0:6,samples=100]{20*x-5*x^2};
\addplot[only marks,mark=*,mark size=1.5pt,DeepBlue] coordinates{(0,0)(.5,8.75)(1,15)(1.5,18.75)(2,20)(2.5,18.75)(3,15)(3.5,8.75)(4,0)(4.5,-11.25)(5,-25)(5.5,-41.25)(6,-60)};
\addplot[guide] coordinates{(2,-65)(2,20)};
\node[anchor=south west,DeepBlue] at (axis cs:2,20) {最高点};
\end{axis}\end{tikzpicture}\end{center}
\end{xesolution}

\begin{origexamplebox}{三种运动函数的图像}\label{ch02:ex:k02-2-2}
在坐标系中分别画出
\[
  (1)\ x=5t^2;\qquad (2)\ x=\frac{3}{t+1};\qquad
  (3)\ x=5\sin(2\pi t).
\]
\par 题图见图~\ref{ch02:ref:question-figure-2-2}。
\KQuestionAC
\end{origexamplebox}
\Needspace{6\baselineskip}\begin{xesolution}
作函数图像不能只凭“形状印象”，应先确定定义域、关键点和变化趋势。若 $t$ 表示从计时开始后的时间，物理上通常只取 $t\ge0$；下面同时说明函数在完整数学定义域上的特征。

\textbf{（1）$x=5t^2$。}
它是开口向上的抛物线，顶点为 $(0,0)$。在 $t\ge0$ 上，
\[
 \frac{\dd x}{\dd t}=10t\ge0,
\]
所以位置随时间单调增加；且斜率 $10t$ 越来越大，图线越来越陡。可先取
$(0,0),(0.5,1.25),(1,5)$ 等点，再用光滑曲线连接。

\textbf{（2）$x=3/(t+1)$。}
完整定义域不含 $t=-1$。当 $t\to-1^+$ 时 $x\to+\infty$，当
$t\to-1^-$ 时 $x\to-\infty$，故 $t=-1$ 是竖直渐近线；当
$t\to\pm\infty$ 时 $x\to0$，故 $x=0$ 是水平渐近线。若只画物理区间
$t\ge0$，则从 $(0,3)$ 出发，并因
\[
 \frac{\dd x}{\dd t}=-\frac{3}{(t+1)^2}<0
\]
而单调下降，逐渐逼近但不会到达 $x=0$。

\textbf{（3）$x=5\sin(2\pi t)$。}
振幅为 $5$，角频率为 $2\pi\,\mathrm{rad/s}$，所以周期
\[
 T=\frac{2\pi}{2\pi\,\mathrm{rad/s}}=1\,\mathrm s.
\]
一个周期内的关键点为
\[
 (0,0),\quad (0.25\,\mathrm s,5),\quad (0.50\,\mathrm s,0),
 \quad (0.75\,\mathrm s,-5),\quad (1.00\,\mathrm s,0).
\]
按这些点画出正弦曲线，再以 $1\,\mathrm s$ 为周期重复。三幅图的纵轴都是位置 $x$，曲线形状不是质点在空间中的轨迹。
\begin{center}\begin{tikzpicture}
\foreach \dx/\mode in {-4.6/1,0/2,4.6/3}{
\begin{scope}[xshift=\dx cm]
\draw[kaxis] (-1.7,0)--(1.75,0) node[right,clearlabel]{$t$};
\draw[kaxis] (-1.45,-1.35)--(-1.45,1.5) node[above,clearlabel]{$x$};
\ifnum\mode=1 \draw[kcurve] plot[domain=0:1.55,samples=60] ({\x-1.45},{.45*\x*\x});\node[clearlabel,below] at (0,-1.6){$5t^2$};\fi
\ifnum\mode=2 \draw[kcurve] plot[domain=0:3.1,samples=80] ({\x-1.45},{1.15/(\x+1)});\node[clearlabel,below] at (0,-1.6){$3/(t+1)\ (t\ge0)$};\fi
\ifnum\mode=3 \draw[kcurve] plot[domain=0:3.1,samples=100] ({\x-1.45},{.85*sin(360*\x)});\node[clearlabel,below] at (0,-1.6){$5\sin2\pi t$};\fi
\end{scope}}
\end{tikzpicture}\end{center}
\end{xesolution}

\begin{origexamplebox}{由位置图判断运动特征}\label{ch02:ex:k02-2-3}
下图分别给出两个质点的一维位置随时间的变化。指出各自运动最明显的特征。
\par 题图见图~\ref{ch02:ref:question-figure-2-3}。
\KQuestionAD
\end{origexamplebox}
\Needspace{6\baselineskip}\begin{xesolution}
先明确读图规则：纵坐标给位置 $x$，曲线斜率给速度 $v=\dd x/\dd t$，极大值和极小值之间的竖直距离不能直接当成路程。

\textbf{第一图。}相邻同向过零点之间的时间间隔不变，说明周期 $T$ 不变；正、负极值到平衡位置 $x=0$ 的距离相同且不随时间改变，说明振幅不变。每到极值点，切线水平，故 $v=0$，质点在这一瞬间换向；经过 $x=0$ 附近时图线最陡，故速度大小最大。一个周期内，质点从一侧极值出发，经过平衡位置，到达另一侧极值，再返回原位置。

\textbf{第二图。}质点仍反复经过 $x=0$ 并在正、负两侧换向，但各次极值的绝对值逐渐减小，所以振幅随时间衰减。若相邻同向过零点的间隔近似相同，只能说周期近似不变；仅凭示意图不能读出精确周期。极值越来越接近零，表示运动逐渐趋向平衡位置，而不是“波形沿水平方向传播”。

两幅图都是同一质点的 $x$--$t$ 图：横向代表时间流逝，不能把图线误认为质点在空间中走出的波浪形轨迹。若坐标轴没有给比例，还只能判断上述定性特征，不能计算具体速度数值。
题图~\ref{ch02:ref:question-figure-2-3}中，左图振幅保持不变；右图各次极值逐渐靠近平衡位置。
\end{xesolution}

上述方法也能描写沿弯曲公路行驶的汽车。只需在公路上选定里程原点和前进的正方向，用沿公路量得的里程标记位置；我国公路上的里程标志就是这种一维坐标思想的实际应用。此时坐标沿道路累计，并不等于汽车到原点的直线距离。

\subsection{速率与速度}

路程 $s$ 是质点沿轨迹实际走过的长度。匀速运动中，速率可写成 $v=s/t$。一般运动在时间区间 $[t,t+\Delta t]$ 内的平均速率为
\begin{equation}
  \bar v_{\mathrm{speed}}=\frac{\Delta s}{\Delta t},
\end{equation}
瞬时速率是时间间隔趋于零时的极限：
\begin{equation}
  v_{\mathrm{speed}}(t)=\lim_{\Delta t\to0}\frac{\Delta s}{\Delta t}.
\end{equation}
在极短时间内，路程可近似为 $\Delta s\simeq v_{\mathrm{speed}}\Delta t$；时间间隔越小，忽略的高阶项越小。

速率是运动快慢的量度，可以把汽车速率表的读数看成直观例子。对匀速运动，用总路程除以总时间即可得到任何时段相同的速率；对变速运动，同一个过程取不同时间段会得到不同的平均速率，所以必须进一步引入瞬时速率。以 $x=5t^2$ 且质点在所考察时段始终沿正方向运动为例，从 $t=1\,\mathrm{s}$ 起选取不同的 $\Delta t$，有
\[
\begin{array}{c|ccccc}
\Delta t/\mathrm s&1&0.1&0.01&0.001&10^{-n}\\\toprule
\Delta s/\Delta t\,/\mathrm{(m/s)}&15&10.5&10.05&10.005&10+5\times10^{-n}
\end{array}
\]
当 $\Delta t\to0$ 时，这些平均值趋于 $10\,\mathrm{m/s}$，即 $t=1\,\mathrm{s}$ 的瞬时速率。这个极限并不是说实际测量要取“等于零的时间”，而是说明时间窗越小，区间平均量越逼近该时刻的运动状态。

同一例子也说明 $\Delta s\simeq v\Delta t$ 的近似。取 $t=1\,\mathrm{s}$ 时的瞬时速率 $10\,\mathrm{m/s}$，当 $\Delta t=1,0.1,0.01\,\mathrm{s}$ 时，线性估计的路程依次为 $10,1,0.1\,\mathrm{m}$，真实路程依次为 $15,1.05,0.1005\,\mathrm{m}$；时间窗缩小后，两者差别迅速减小。数学上这正是切线对曲线的局部近似。

一维坐标轴已经规定了正方向。为了同时表达快慢和方向，引入有正负号的速度
\begin{equation}
  \bar v=\frac{\Delta x}{\Delta t},\qquad
  v(t)=\lim_{\Delta t\to0}\frac{\Delta x}{\Delta t}=\frac{\dd x}{\dd t}.
\end{equation}
瞬时速率等于 $|v|$。若速度在区间内恒定，则
\begin{equation}
  x=x_0+v\Delta t.
\end{equation}

例如 $x=5t$ 的平均速度与所取区间无关，始终为 $5\,\mathrm{m/s}$；而 $x=5t^2$ 在 $t=1\,\mathrm{s}$ 附近的平均速度为 $10+5\Delta t$，当 $\Delta t\to0$ 时趋于 $10\,\mathrm{m/s}$。这正是从平均量过渡到瞬时量的极限思想。

必须区分路程、位移、速率和速度。路程记录轨迹累计长度，始终非负；位移是末位置减初位置，在一维中可正可负。速率只表示快慢，速度还包含方向。若质点先沿正方向走一段再返回，路程继续增加，而位移可能减小甚至回到零；相应地，平均速率一般不等于平均速度的绝对值。式 $x=x_0+v\Delta t$ 对正、负速度都成立，符号已经把方向包含在内。

例如质点从 $x=0$ 走到 $x=10\,\mathrm m$，再返回 $x=4\,\mathrm m$，全过程路程为 $16\,\mathrm m$，位移却只有 $+4\,\mathrm m$。若用时 $4\,\mathrm s$，平均速率是 $4\,\mathrm{m/s}$，平均速度是 $+1\,\mathrm{m/s}$。只有运动过程中不换向时，路程才等于位移大小，平均速率才等于平均速度的绝对值。判断是否换向应看速度符号，而不能只看位置始终为正还是为负。

瞬时速度为零也不等于质点将在此后保持静止。位置图像在极值点切线水平，表示该时刻速度为零；若曲线随后斜率改变符号，质点正是在这一瞬间换向。相反，若位置图在一段时间内保持水平，才表示质点在这一整段都静止。这一区别在竖直上抛最高点和往复振动中尤其重要。

\begin{origexamplebox}{相邻测量间的平均速率}\label{ch02:ex:k02-2-4}
表中 $t=0,0.5,\ldots,6\,\mathrm{s}$ 时的位置分别为
\[
0,8.75,15,18.75,20,18.75,15,8.75,0,-11.25,-25,-41.25,-60\ \mathrm{m}.
\]
求相邻两次测量之间的平均速率。
\end{origexamplebox}
\Needspace{6\baselineskip}\begin{xesolution}
相邻两次测量的时间间隔均为
\[
 \Delta t=0.5\,\mathrm s.
\]
先计算每一小段的位移 $\Delta x=x_{i+1}-x_i$，再除以 $\Delta t$。例如前两段为
\[
 \bar v_{0\sim0.5}=\frac{8.75-0}{0.5}=17.5\,\mathrm{m/s},
 \qquad
 \bar v_{0.5\sim1}=\frac{15-8.75}{0.5}=12.5\,\mathrm{m/s}.
\]
依次计算得到有符号平均速度
\[
17.5,12.5,7.5,2.5,-2.5,-7.5,-12.5,-17.5,-22.5,-27.5,-32.5,-37.5\ \mathrm{m/s}.
\]
由例~\ref{ch02:ex:k02-2-1}的抛物线可知，物体在每个相邻测量区间内都不发生额外换向，唯一的换向点恰在区间端点 $t=2\,\mathrm s$。因此各段路程等于 $|\Delta x|$，平均速率就是上述数值的绝对值：
\[
\boxed{17.5,12.5,7.5,2.5,2.5,7.5,12.5,17.5,22.5,27.5,32.5,37.5\ \mathrm{m/s}}.
\]
$t=2\,\mathrm{s}$ 前平均速度为正，之后为负；负号表示沿坐标轴负方向运动，不表示运动“比静止更慢”。还要注意：若只知道相邻时刻的位置，却不知道区间内是否换向，一般只能求平均速度，不能仅凭端点确定平均速率；本题能这样做，是因为前一例已经确定了区间内的运动模型。
\end{xesolution}

\begin{origexamplebox}{相向运动}\label{ch02:ex:k02-2-5}
质点 $A,B$ 在同一直线上匀速相向运动。初始位置 $x_A=-5\,\mathrm m$、$x_B=10\,\mathrm m$；$A$ 沿正方向以 $1.5\,\mathrm{m/s}$ 运动，$B$ 沿负方向的速率为 $0.5\,\mathrm{m/s}$。求相遇时间和位置。
\end{origexamplebox}
\Needspace{6\baselineskip}\begin{xesolution}
取题目给定的正方向，令 $t=0$ 为两质点处在初始位置的时刻。匀速直线运动满足“位置等于初位置加上速度乘时间”，因此
\[
x_A(t)=-5\,\mathrm m+(1.5\,\mathrm{m/s})t,
\qquad
x_B(t)=10\,\mathrm m-(0.5\,\mathrm{m/s})t.
\]
相遇的物理条件是同一时刻坐标相同：
\[
-5+1.5t=10-0.5t.
\]
整理得 $(2.0\,\mathrm{m/s})t=15\,\mathrm m$，所以
\[
\boxed{t=7.5\,\mathrm s},\qquad
\boxed{x=6.25\,\mathrm m}.
\]
其中相遇位置可分别代回两式检查：
\[
x_A=-5+1.5\times7.5=6.25\,\mathrm m,
\qquad
x_B=10-0.5\times7.5=6.25\,\mathrm m.
\]
再用相对运动复核：初始间距为 $15\,\mathrm m$，接近速率为
$1.5+0.5=2.0\,\mathrm{m/s}$，故闭合时间为 $15/2.0=7.5\,\mathrm s$，与坐标法一致。时间量纲也由“长度除以速度”正确得到。
\end{xesolution}

\subsection{图像的斜率与面积}

\PositionSlopeDiagram

$x$-$t$ 图像中两点割线斜率是平均速度，切线斜率是瞬时速度。图线向上并不一定表示质点在空间中“向上”，它只表示坐标随时间增大；具体方向由坐标轴的正方向决定。

图~\ref{ch02:fig:k02-2-3}把这一点作了几何化说明。匀速运动的位置图像是一条直线，任意两点割线都与原直线重合，斜率恒定。一般运动中，时刻 $t$ 与 $t+\Delta t$ 对应曲线上的 $A,B$ 两点，割线 $AB$ 的斜率是 $\Delta x/\Delta t$。让 $\Delta t$ 逐渐缩小，$B$ 沿曲线趋近 $A$，割线趋近于 $A$ 点切线，其斜率的极限就是瞬时速度。因而位置图像的几何信息可以直接翻译成速度：斜率为正表示沿正方向，斜率为负表示沿负方向，斜率为零表示该时刻瞬时静止。

\begin{origexamplebox}{从 $x$-$t$ 图判断速度趋势}\label{ch02:ex:k02-2-6}
根据下列三幅位置—时间图像，分别判断质点的运动方向和速度大小怎样变化。
\par 题图见图~\ref{ch02:ref:question-figure-2-6}。
\KQuestionAG
\end{origexamplebox}
\Needspace{6\baselineskip}\begin{xesolution}
沿横轴箭头方向读取时间，并在每一段比较切线斜率 $v=\dd x/\dd t$。


\ExampleTwoSixGraphs
\textbf{第一图：先线性上升后线性下降。}上升直线的斜率为正常量，表示质点先以恒定正速度运动；下降直线的斜率为负常量，表示随后以恒定负速度运动。折点处左右斜率不同，理想模型中速度在同一瞬间跳变，普通加速度不存在；真实换向若由有限力造成，折角附近应是一小段光滑曲线。

\textbf{第二图：向上凹且越来越陡。}任一点切线斜率均为正，所以质点始终沿正方向运动；斜率随时间增大，所以正速度增大，因而加速度为正。这里“图线向上弯”是对 $x$--$t$ 曲线凹凸性的描述，不代表质点在空间中沿弯曲轨迹运动。

\textbf{第三图：上半圆形。}左半段斜率为正，但逐渐减小到零，表示质点沿正方向运动并逐渐减速；最高点切线水平，故该时刻 $v=0$；右半段斜率为负且绝对值逐渐增大，表示质点换向后沿负方向越来越快。若把理想半圆写成
\[
 \left(\frac{t-t_c}{\tau}\right)^2+
 \left(\frac{x-x_c}{X}\right)^2=1,\qquad x\ge x_c,
\]
则
\[
 v=\frac{\dd x}{\dd t}
   =-\frac X\tau\frac{(t-t_c)/\tau}{\sqrt{1-[(t-t_c)/\tau]^2}}.
\]
这里 $\tau>0$ 是时间尺度，$X>0$ 是长度尺度，不能把时间平方与位置平方直接相加。
图中“半圆”指归一化坐标 $((t-t_c)/\tau,(x-x_c)/X)$ 下的几何形状；
实际位置为 $x=x_c+X\sqrt{1-[(t-t_c)/\tau]^2}$，求导后含有速度尺度 $X/\tau$。
这验证了顶点 $t=t_c$ 时 $v=0$，而两端 $t=t_c\pm\tau$ 的竖直切线处 $|v|\to\infty$。
无限速度没有现实意义，说明把一整段真实位置图精确画成半圆通常只是一种几何理想化。
\end{xesolution}

速度也可以画成 $v$-$t$ 图。时间区间内的位移为
\begin{equation}
  \Delta x=\int_{t_1}^{t_2}v(t)\,\dd t,
\end{equation}
即曲线与时间轴围成的\textbf{有向面积}；路程则为 $\int|v|\dd t$。

图~\ref{ch02:fig:k02-2-4}还说明了面积结论的来由。图中 $A,B$ 表示时间区间的两个端点，$C,D$ 是相应时刻速度曲线上的点；因此封闭区域 $ABDC$ 正是这一时间段的累计位移。匀速时，$v$-$t$ 图是一条平行于时间轴的直线，矩形面积 $v\Delta t$ 就是位移。变速时，把时间区间分成许多很窄的小段，每段内用左端或某个时刻的瞬时速度近似不变，小矩形面积近似该小段位移。分割越细，所有小矩形面积之和越接近曲线与时间轴围成的面积，极限就是积分。曲线位于时间轴下方时，速度为负，对位移贡献负面积；计算总路程则要把这部分面积取正。

\VelocityAreaDiagram

\begin{origexamplebox}{速度图与位移图}\label{ch02:ex:k02-2-7}
速度从正值随时间线性减小，在 $T/2$ 时变为零，$T$ 时变为等大的负值。判断从 $x(0)=0$ 开始的位置图像。
\par 题图见图~\ref{ch02:ref:question-figure-2-7}。
\KQuestionAH
\end{origexamplebox}
\Needspace{6\baselineskip}\begin{xesolution}
设初始速度为 $v_0>0$，末速度为 $-v_0$。速度随时间线性变化，所以连接 $(0,v_0)$ 与 $(T,-v_0)$ 的直线方程是
\[
 v(t)=v_0-\frac{2v_0}{T}t.
\]
\ExampleTwoSevenGraphs

其斜率即加速度，
\[
 a=-\frac{2v_0}{T}<0,
\]
为常量。因此位置由积分得到
\[
 x(t)-x(0)=\int_0^t v(t')\,\dd t'
 =v_0t-\frac{v_0}{T}t^2.
\]
又因 $x(0)=0$，
\[
 \boxed{x(t)=v_0t-\frac{v_0}{T}t^2},
\]
这是开口向下的抛物线。它在 $t=T/2$ 处斜率为零并达到最大位置
\[
 x_{\max}=x(T/2)=\frac{v_0T}{4},
\]
在 $t=T$ 时回到 $x(T)=0$。

同一结论也可由 $v$--$t$ 图的有向面积看出：$0$ 到 $T/2$ 的正三角形面积与 $T/2$ 到 $T$ 的负三角形面积大小相等，故总位移为零。右侧四图中应选光滑的抛物线拱形 B 图；折线 A、D 对应速度突变，C 图的末位置不回到零。式中 $v_0t$ 与 $(v_0/T)t^2$ 都是长度，量纲一致。
\end{xesolution}

\begin{origexamplebox}{由 $v$-$t$ 图求路程}\label{ch02:ex:k02-2-8}
求三种运动在 $0$ 到 $t$ 内的路程：（1）速度从 $0$ 线性增至 $v$；（2）速度从 $v_0$ 线性增至 $v_1$；（3）速度从 $0$ 线性增至 $v_0$，在 $t/3$ 达到峰值，再线性降到 $0$。
\par 题图见图~\ref{ch02:ref:question-figure-2-8}。
\KQuestionAI
\end{origexamplebox}
\Needspace{6\baselineskip}\begin{xesolution}
三幅 $v$--$t$ 图都位于时间轴上方，故速度始终非负，路程与位移相等，均可直接用图线下的面积计算。


\ExampleTwoEightGraphs
\textbf{（1）}图形是底为 $t$、高为 $v$ 的三角形，
\[
 s_1=\frac12tv.
\]

\textbf{（2）}图形是两条平行边长分别为 $v_0,v_1$、高为 $t$ 的梯形，
\[
 s_2=\frac{v_0+v_1}{2}t.
\]

\textbf{（3）}峰值左、右两部分分别是底为 $t/3$ 与 $2t/3$、共同高为 $v_0$ 的三角形，
\[
 s_3=\frac12v_0\frac t3+\frac12v_0\frac{2t}{3}
 =\frac12v_0t.
\]
综上，
\[
  \boxed{s_1=\frac12vt},\qquad
  \boxed{s_2=\frac{v_0+v_1}{2}t},\qquad
  \boxed{s_3=\frac12v_0t}.
\]
第三问说明：只要速度从零线性升到同一峰值再线性降到零，峰值时刻只是在左右两个三角形之间重新分配底边，总底边仍为 $t$，所以总面积不变。各答案都是“速度乘时间”，量纲为长度。若某段图线落在时间轴下方，求位移要保留负面积，求路程则必须把该部分面积取正；本题没有这种情况。
\end{xesolution}

\subsection{加速度与匀变速直线运动}

速度变化量为 $\Delta v=v(t+\Delta t)-v(t)$。平均加速度和瞬时加速度分别为
\begin{equation}
  \bar a=\frac{\Delta v}{\Delta t},\qquad
  a(t)=\lim_{\Delta t\to0}\frac{\Delta v}{\Delta t}=\frac{\dd v}{\dd t}.
\end{equation}
加速度描述速度变化的快慢和方向。$a$ 与 $v$ 同号时速率增加，异号时速率减小；不能把“加速度为负”直接说成“减速”。

与速度的定义完全类似，若 $\Delta t$ 取一个有限区间，$\Delta v/\Delta t$ 是该区间的平均加速度；若让时间间隔趋于零，就得到某一时刻的瞬时加速度。相同时间内速度变化越大，加速度的绝对值越大。加速度的正负由坐标正方向决定，它表示速度在数轴上的变化方向，而不是单独表示“加快”或“减慢”。例如 $v<0,a<0$ 时，速度的绝对值增大，质点实际在加速向负方向运动。

速度图像为正只说明运动方向沿坐标正向，为负只说明沿负向；加速度图像同理表示速度的变化方向。判断速率增减，可以比较 $v$ 与 $a$ 的符号：同号时 $|v|$ 增大，异号时 $|v|$ 减小。若 $v=0$ 而 $a\ne0$，质点可能正处在换向的临界时刻。把这些符号关系与 $x$-$t$、$v$-$t$、$a$-$t$ 三类图像同时核对，能排除许多仅凭“图线向上或向下”产生的误解。

三个运动学量之间存在成对的局部与累计关系：位置图的切线斜率给速度，速度图的切线斜率给加速度；速度图的有向面积给位移，加速度图的有向面积给速度变化量。所谓有向面积必须保留横轴上下的正负。若只求路程，则须把速度负值部分翻到横轴上方后再累计，而不能在最后对总位移简单取绝对值。

实际测量的“瞬时值”总是在很短但有限的时间窗内取得的近似。速度表、光电门和高速摄影并不会真的用零秒作分母，而是让采样时间足够短，使这段时间内速度变化小于所要求的精度。数学极限给出理想定义，实验误差则决定时间窗能缩小到什么程度；时间窗过大不能代表瞬时状态，过小又可能放大计时和位置测量误差。

从函数观点看，$x(t)$ 可导时，$v(t)$ 就是其导函数；$v(t)$ 可导时，$a(t)$ 是速度的导函数。因此位置连续不保证速度连续，速度连续也不保证加速度连续。理想碰撞或分段控制常使速度图出现折点甚至跳变，这时某一瞬间的普通加速度可能不存在，需要按碰撞前后分段讨论。讲义中的匀变速公式默认所考察时间段内加速度有限且恒定。

消去时间得到的关系
\begin{equation}
 v^2-v_0^2=2a(x-x_0)
\end{equation}
特别适合题目不给时间的情形。它含有速度平方，求最终速度时仍要结合运动方向选择正、负根；只写 $v=\sqrt{v_0^2+2a\Delta x}$ 会默认最终沿正方向。若物体先减速到零再反向，同一段过程还应以换向时刻为界分段，避免把一条带符号的位移误当成全程路程。

匀变速运动中，任意时间段的平均速度等于首末速度的算术平均，但这个结论依赖加速度恒定。一般变速运动不能只用 $(v_1+v_2)/2$ 代替区间平均速度，应回到 $v$-$t$ 图面积或积分。相应地，列车中点通过定点的速度不是首尾速度的算术平均，正是因为题目给的是相等位移而非相等时间。

列运动方程时还应先统一单位。若距离用米、时间用秒，速度和加速度应分别化为 $\mathrm{m/s}$ 与 $\mathrm{m/s^2}$；千米每小时不能直接与米、秒混算。答案完成后，用量纲检查每一项：位置方程各项必须都是长度，速度方程各项必须都是速度。量纲正确不能保证数值一定正确，却能迅速发现漏乘时间、误用平方或单位换算错误。

若 $a$ 恒定，初始位置和速度为 $x_0,v_0$，则
\begin{equation}
  v=v_0+at,\qquad
  x=x_0+v_0t+\frac12at^2.
\end{equation}
$v$-$t$ 图的斜率是 $a$，$a$-$t$ 图的面积是速度变化量。

匀加速运动的 $v$-$t$ 图是直线，斜率为常量 $a$；直线下的梯形面积给出位移，得到位置的二次函数。一般运动的 $v$-$t$ 曲线在某点的切线斜率给出该时刻加速度；反过来，给出 $a$-$t$ 图后，两时刻间曲线的有向面积等于速度改变量 $v_2-v_1$。这与“位置图的斜率给速度、速度图的面积给位移”组成一组相互对应的微分与积分关系。

\begin{origexamplebox}{用速度图推出位置公式}\label{ch02:ex:k02-2-9}
根据匀加速运动的速度-时间图像，由 $v=v_0+at$ 推出位置随时间的关系。
\end{origexamplebox}
\Needspace{6\baselineskip}\begin{xesolution}
匀加速运动满足 $v(t)=v_0+at$，所以 $v$--$t$ 图是从 $(0,v_0)$ 到
$(t,v)$ 的直线。位置变化等于速度图下的有向面积：
\[
 x-x_0=\int_0^t v(t')\,\dd t'.
\]
几何上，这一面积是一个梯形；也可拆成高为 $v_0$ 的矩形和高为 $at$ 的三角形：
\[
 x-x_0=v_0t+\frac12(at)t
 =v_0t+\frac12at^2.
\]
若直接用梯形面积，同样有
\[
  x-x_0=\frac{v_0+v}{2}t
  =\frac{v_0+(v_0+at)}{2}t
  =v_0t+\frac12at^2.
\]
因此
\[
 \boxed{x=x_0+v_0t+\frac12at^2}.
\]
检查：对右式求时间导数，得到 $\dd x/\dd t=v_0+at=v$，与原速度式一致；令 $a=0$，又退化为匀速公式 $x=x_0+v_0t$。量纲方面，$v_0t$ 与 $at^2$ 都是长度。图线若部分位于时间轴下方，仍须按有向面积计算，公式中的速度和加速度都要保留符号。
\end{xesolution}

\begin{origexamplebox}{消去时间}\label{ch02:ex:k02-2-10}
匀加速质点由 $(x_1,v_1)$ 运动到 $(x_2,v_2)$，证明
\[
  v_2^2-v_1^2=2a(x_2-x_1).
\]
\end{origexamplebox}
\Needspace{6\baselineskip}\begin{xesolution}
记从状态1到状态2所经历的时间为 $\Delta t$。由于加速度恒定，
\[
 v_2-v_1=a\Delta t,
\]
而这段时间内的速度图是一条直线，因此平均速度是首、末速度的算术平均：
\[
  x_2-x_1=\frac{v_1+v_2}{2}\Delta t.
\]
将第二式乘以 $2a$，再用 $a\Delta t=v_2-v_1$ 消去时间：
\[
\begin{aligned}
2a(x_2-x_1)
 &=a(v_1+v_2)\Delta t\\
 &=(v_2-v_1)(v_2+v_1)\\
 &=v_2^2-v_1^2.
\end{aligned}
\]
故
\[
 \boxed{v_2^2-v_1^2=2a(x_2-x_1)}.
\]
两边量纲都是 $\mathrm{m^2/s^2}$。该关系不含时间，适用于所考察整段加速度恒定的情形；若由它反求 $v_2$，必须结合运动方向选择平方根的正负号，不能总写成正根。
\end{xesolution}

\begin{origexamplebox}{列车中部通过定点的速度}\label{ch02:ex:k02-2-11}
一列火车以恒定加速度沿直线前进。车头通过点 $O$ 时速度为 $v_1$，车尾通过 $O$ 时速度为 $v_2$。求列车中部通过 $O$ 时的速度 $v$。
\end{origexamplebox}
\Needspace{6\baselineskip}\begin{xesolution}
设车长为 $L$，取列车前进方向为正。把整列车看成刚体，用车头、车中部或车尾上任一点的位置描述列车平移，三点在同一时刻具有相同速度。

车头刚通过 $O$ 时，把该时刻记为状态1。等车尾通过 $O$ 时，整列车恰向前平移了一个车长 $L$，于是由无时间公式
\[
  v_2^2-v_1^2=2aL.
\]
当列车中部通过 $O$ 时，列车相对状态1只平移了 $L/2$，故
\[
 v^2-v_1^2=2a\frac L2=aL=\frac{v_2^2-v_1^2}{2}.
\]
整理得
\[
 v^2=\frac{v_1^2+v_2^2}{2}.
\]
列车始终沿正方向前进，所以取正根：
\[
  \boxed{v=\sqrt{\frac{v_1^2+v_2^2}{2}}}.
\]
这不是 $(v_1+v_2)/2$，因为中部对应的是总位移的一半，不是总时间的一半。检查：若 $a=0$，则 $v_1=v_2$，上式正确退化为 $v=v_1$；若 $v_2>v_1$，所得 $v$ 确在 $v_1$ 与 $v_2$ 之间。
\end{xesolution}

\begin{origexamplebox}{通过路口或制动停车}\label{ch02:ex:k02-2-12}
汽车以 $v_0=20\,\mathrm{m/s}$ 行驶，在距路口 $60\,\mathrm m$ 处发现倒计时还剩 $2\,\mathrm s$。忽略司机反应时间，并假定采取方案后加速度保持恒定。若加速通过，所需加速度至少多大？若停车不闯红灯，制动加速度大小至少多大？
\end{origexamplebox}
\Needspace{6\baselineskip}\begin{xesolution}
两种选择要分别建立临界条件。“至少”的边界都对应恰好满足要求；比临界值更小就会失败。

\textbf{加速通过。}设加速度沿行驶方向为 $a_{\rm pass}>0$。要在倒计时结束时刚好到达路口，$2\,\mathrm s$ 内的位移应恰为 $60\,\mathrm m$：
\[
60\,\mathrm m=(20\,\mathrm{m/s})(2\,\mathrm s)
+\frac12a_{\rm pass}(2\,\mathrm s)^2.
\]
所以
\[
\boxed{a_{\rm pass}=10\,\mathrm{m/s^2}}.
\]
此时末速度为 $20+10\times2=40\,\mathrm{m/s}>0$，运动方向没有改变。

\textbf{制动停车。}设制动加速度的大小为 $a_{\rm stop}>0$，实际加速度为 $-a_{\rm stop}$。最小制动能力对应汽车恰在路口前速度降为零、制动距离恰为 $60\,\mathrm m$。使用无时间公式：
\[
0^2-(20\,\mathrm{m/s})^2
=2(-a_{\rm stop})(60\,\mathrm m),
\]
故
\[
\boxed{a_{\rm stop}=\frac{10}{3}\,\mathrm{m/s^2}}.
\]
复核制动时间为 $t=v_0/a_{\rm stop}=6\,\mathrm s$；停车方案不要求在剩余 $2\,\mathrm s$ 内停下，只要求不越过路口。两种计算都忽略了司机反应时间：若反应阶段仍以 $20\,\mathrm{m/s}$ 前进，剩余制动距离会缩短，实际所需制动加速度将更大。
\end{xesolution}

\begin{origexamplebox}{追尾避免条件}\label{ch02:ex:k02-2-13}
$A$ 车以 $20\,\mathrm{m/s}$ 行驶，发现前方同轨道上相距 $100\,\mathrm m$ 的另一列车以 $10\,\mathrm{m/s}$ 行驶，立即以大小为 $a$ 的恒定加速度制动。求不相撞所需的最小 $a$。
\end{origexamplebox}
\Needspace{6\baselineskip}\begin{xesolution}
取行驶方向为正，设后车为 $A$、前车为 $B$。从开始制动计时，
\[
x_A=20t-\frac12at^2,
\qquad
x_B=100+10t.
\]
两车间距为“前车位置减后车位置”：
\[
d(t)=x_B-x_A=100-10t+\frac12at^2.
\]
只要 $d(t)\ge0$，后车就没有越过前车。对时间求导，
\[
\frac{\dd d}{\dd t}=-10+at.
\]
所以间距在
\[
t_*=\frac{10}{a}
\]
时最小；这也正是两车速度相等、相对速度降为零的时刻。代回得
\[
d_{\min}=100-10\frac{10}{a}
+\frac12a\left(\frac{10}{a}\right)^2
=100-\frac{50}{a}.
\]
不相撞要求最小间距非负：
\[
100-\frac{50}{a}\ge0,\qquad \boxed{a_{\min}=0.50\,\mathrm{m/s^2}}.
\]
也可用相对运动的“制动距离”快速复核：初始接近速度为
$20-10=10\,\mathrm{m/s}$，相对减速度大小为 $a$，消除相对速度所需的闭合距离为
\[
\frac{(10\,\mathrm{m/s})^2}{2a}\le100\,\mathrm m,
\]
得到同一条件。临界情况下 $a=0.50\,\mathrm{m/s^2}$、$t_*=20\,\mathrm s$，两车以共同速度 $10\,\mathrm{m/s}$ 恰好相接；若要求保留正的安全间距，则应取 $a>0.50\,\mathrm{m/s^2}$。
\end{xesolution}

\begin{origexamplebox}{直道追及条件}\label{ch02:ex:k02-2-14}
两辆赛车进入长度为 $L$ 的直道，并在所考察时间内始终沿直道正方向运动。以两车车头所在点代表各自位置：初始时后车位于直道起点，前车在其前方 $d$ 处；后车初速为 $v_1$，前车初速为 $v_2>v_1$；后车加速度为 $a_1>0$，前车加速度为 $a_2$，且 $a_1>a_2$。求后车车头在到达直道末端以前追上前车车头的条件。若要求车身完全超越，还需另给两车车长，本题不讨论该情形。
\end{origexamplebox}
\Needspace{6\baselineskip}\begin{xesolution}
取后车进入直道的位置为 $x=0$，令前车初始位置为 $x=d$。两车的位置分别为
\[
x_1(t)=v_1t+\frac12a_1t^2,
\qquad
x_2(t)=d+v_2t+\frac12a_2t^2.
\]
设
\[
\Delta v=v_2-v_1>0,
\qquad
\Delta a=a_1-a_2>0.
\]
前车与后车的间距是
\[
 r(t)=x_2-x_1=d+\Delta v\,t-\frac12\Delta a\,t^2.
\]
开始时 $r(0)=d>0$，且前车初速度较大，间距会先增大；但后车加速度更大，最终有可能追上。相遇条件为 $r(t)=0$：
\[
\frac12\Delta a\,t^2-\Delta v\,t-d=0.
\]
两个根中只有下式为正，故第一次相遇时刻为
\[
 t_*=\frac{\Delta v+\sqrt{\Delta v^2+2\Delta a d}}{\Delta a}.
\]
相遇点坐标等于后车走过的距离
\[
x_*=x_1(t_*)=v_1t_*+\frac12a_1t_*^2.
\]
因此，后车能在直道末端以前追上前车的充要条件是
\[
  \boxed{v_1t_*+\frac12a_1t_*^2<L}.
\]
另一种检验方式是不先求相遇点。令后车到达直道末端的时间为 $t_L$，即
\[
L=v_1t_L+\frac12a_1t_L^2.
\]
因题设 $a_1>0$，取正根得
\[
t_L=\frac{-v_1+\sqrt{v_1^2+2a_1L}}{a_1},
\]
则条件等价为
\[
\boxed{r(t_L)=d+\Delta v\,t_L-\frac12\Delta a\,t_L^2<0}.
\]
因为 $r(t)$ 的正根只有一个，$r(t_L)<0$ 恰好表示相遇已经在 $t_L$ 之前发生；若 $r(t_L)=0$，只能说明两车车头恰在直道末端重合，不满足“末端以前”。量纲检查中，$\Delta vt$、$\Delta at^2$、$d$ 和 $L$ 都是长度。若计算得到 $x_*\ge L$，即使无限长道路上最终能追上，也不能算在本段直道内完成追及。
\end{xesolution}

\subsection{落体运动}

取竖直向上为正方向，忽略空气阻力时，地表附近所有自由落体具有相同加速度
\begin{equation}
  a=-g,\qquad g\approx9.8\,\mathrm{m/s^2}.
\end{equation}
若 $t=0$ 时位置和速度为 $x_0,v_0$，则
\begin{equation}
  v=v_0-gt,\qquad x=x_0+v_0t-\frac12gt^2.
\end{equation}
公式中的正负号由坐标正方向决定；若改取竖直向下为正，加速度分量改为 $a=+g$，
相应重力项也取正号。$g$ 始终表示重力加速度的大小，不随坐标反向变成负数。

把物体由静止释放后，在地表附近小范围内可近似认为重力加速度大小不随高度改变，这种运动是匀加速运动的一种。实验表明，在忽略空气阻力时，不同质量和材料的物体具有相同的自由落体加速度。若由原点静止释放，竖直向上为正，则 $v=-gt$、$x=-gt^2/2$；若初始位置或初速度不为零，直接使用上面的通式。所谓“上抛运动”和“下抛运动”并不是新的公式，只是 $v_0$ 符号不同。

空气中羽毛与石块落得不同，主要来自阻力相对重力的影响不同，并不否定真空中自由落体加速度相同。使用恒定 $g$ 的模型还隐含高度变化远小于地球半径；若运动范围扩展到高空，$g$ 的大小和方向都需要随位置修正。在本节范围内，把竖直运动统一写成带符号的匀变速方程，比为“自由落下、上抛、下抛”分别背三套公式更可靠。

\begin{origexamplebox}{三种竖直运动的位置}\label{ch02:ex:k02-2-15}
以抛出点为原点、竖直向上为正，取 $g=10\,\mathrm{m/s^2}$。分别求：（1）静止释放后 $1,5,10\,\mathrm s$ 的位置；（2）以 $10\,\mathrm{m/s}$ 竖直上抛后 $1,2,3\,\mathrm s$ 的位置；（3）以 $10\,\mathrm{m/s}$ 竖直下抛后相同时刻的位置。
\end{origexamplebox}
\Needspace{6\baselineskip}\begin{xesolution}
统一取抛出点为 $x=0$、竖直向上为正。因为 $a=-g=-10\,\mathrm{m/s^2}$，三种运动都使用
\[
x(t)=v_0t-\frac12gt^2=v_0t-5t^2,
\]
区别只在初速度 $v_0$。

\textbf{（1）静止释放：$v_0=0$。}
\[
x=-5t^2.
\]
分别代入 $t=1,5,10\,\mathrm s$，得到
\[
\begin{array}{c|ccc}
t/\mathrm{s}&1&5&10\\\toprule
 x/\mathrm m&-5&-125&-500\\\bottomrule
\end{array}
\]

\textbf{（2）竖直上抛：$v_0=+10\,\mathrm{m/s}$。}
\[
x=10t-5t^2.
\]
在 $t=1,2,3\,\mathrm s$ 时，
\[
\begin{array}{c|ccc}
t/\mathrm{s}&1&2&3\\\toprule
x/\mathrm m&5&0&-15\\\bottomrule
\end{array}
\]
$t=1\,\mathrm s$ 时速度 $v=10-10t=0$，物体处在最高点；$t=2\,\mathrm s$ 时回到原点；随后位置变为负。

\textbf{（3）竖直下抛：$v_0=-10\,\mathrm{m/s}$。}
\[
x=-10t-5t^2,
\]
故在 $t=1,2,3\,\mathrm s$ 时，
\[
\begin{array}{c|ccc}
t/\mathrm{s}&1&2&3\\\toprule
x/\mathrm m&-15&-40&-75\\\bottomrule
\end{array}
\]
负号只说明物体位于原点下方。三种位置公式的每一项都是长度；上抛、下抛无需另背公式，只要把初速度按坐标方向正确取正负即可。
\end{xesolution}

\begin{origexamplebox}{到达指定高度的时间}\label{ch02:ex:k02-2-16}
质点从原点以 $v_0=20\,\mathrm{m/s}$ 竖直上抛，取 $g=10\,\mathrm{m/s^2}$。求到达 $x=15,20,25,0,-25\,\mathrm m$ 所需时间。
\end{origexamplebox}
\Needspace{6\baselineskip}\begin{xesolution}
位置方程为
\[
x=20t-5t^2.
\]
对给定位置 $x$，整理为
\[
5t^2-20t+x=0,
\]
其根为
\[
t=2\pm\sqrt{4-\frac{x}{5}},
\]
并且只保留 $t\ge0$ 的根。

逐项计算：
\[
\begin{aligned}
x=15\,\mathrm m:&\quad t=2\pm1=1,3\,\mathrm s;\\
x=20\,\mathrm m:&\quad t=2\,\mathrm s\quad\text{（重根）};\\
x=25\,\mathrm m:&\quad 4-x/5=-1<0\quad\text{（无实根）};\\
x=0:&\quad t=0,4\,\mathrm s;\\
x=-25\,\mathrm m:&\quad t=2\pm3=-1,5\,\mathrm s,
\end{aligned}
\]
最后一项舍去负时间 $-1\,\mathrm s$。汇总为
\[
\begin{array}{c|ccccc}
x/\mathrm m&15&20&25&0&-25\\\toprule
t/\mathrm s&1,3&2&\text{无实根}&0,4&5\\\bottomrule
\end{array}
\]
为什么 $x=15\,\mathrm m$ 有两个时刻？速度
\[
v=20-10t
\]
在 $t=1\,\mathrm s$ 为正、在 $t=3\,\mathrm s$ 为负，分别对应上升和下降经过同一高度。最高点由 $v=0$ 得 $t=2\,\mathrm s$，最大高度为
\[
x_{\max}=\frac{v_0^2}{2g}=20\,\mathrm m,
\]
所以不可能到达 $25\,\mathrm m$。$x=0$ 的两个时刻分别是抛出瞬间和回到原点时。
\end{xesolution}

\begin{origexamplebox}{一落一上抛}\label{ch02:ex:k02-2-17}
高度 $h$ 处的质点 $A$ 自由下落，同时从地面以初速度 $v_0$ 竖直上抛质点 $B$。允许临界情况下两者恰在 $A$ 落地时相遇，求所需的最小 $v_0$。
\par 题图见图~\ref{ch02:ref:question-figure-2-17}。
\KQuestionAR
\end{origexamplebox}
\Needspace{6\baselineskip}\begin{xesolution}
取地面为 $y=0$、竖直向上为正。两球的位置分别为
\[
y_A=h-\frac12gt^2,
\qquad
y_B=v_0t-\frac12gt^2.
\]
两式中的重力项完全相同，这是因为两球都具有加速度 $-g$。令 $y_A=y_B$，得到
\[
h=v_0t,
\]
所以相遇时刻为
\FreeFallMeetingDiagram
\[
t_m=\frac{h}{v_0}.
\]
$A$ 从高度 $h$ 落到地面的时间由 $h-gt_f^2/2=0$ 得
\[
t_f=\sqrt{\frac{2h}{g}}.
\]
要在 $A$ 落地以前或恰好落地时相遇，必须有 $t_m\le t_f$，即
\[
\frac{h}{v_0}\le\sqrt{\frac{2h}{g}}.
\]
因 $h,v_0>0$，整理得临界值
\[
  \boxed{v_{0,\min}=\sqrt{\frac{gh}{2}}}.
\]
也可从相对运动理解：两球相对加速度为零，初始间距 $h$ 以恒定相对速度 $v_0$ 闭合。临界情况下
$t_m=t_f$，两球恰在地面相遇，所以题设所允许的最小值能够取到。若另加“必须在离地面高度大于零处相遇”的严格条件，则必须取
$v_0>\sqrt{gh/2}$，此时不存在能取到的最小值，上式只是允许初速度的下确界。答案的量纲为
$\sqrt{(\mathrm{m/s^2})\mathrm m}=\mathrm{m/s}$。
\end{xesolution}

\begin{origexamplebox}{超人追赶自由落体}\label{ch02:ex:k02-2-18}
一名女子从高为 $H$ 的楼顶落下。经过时间 $T$ 后，超人从楼顶以向下初速度 $v_0$ 出发，并同样以重力加速度向下运动。允许临界情况下恰在女子落地时追上她，求所需的最小 $v_0$。
\end{origexamplebox}
\Needspace{6\baselineskip}\begin{xesolution}
取楼顶为原点、竖直向下为正，并从女子开始下落时计时。女子的位置为
\[
y_w(t)=\frac12gt^2,
\]
故她的落地时刻为
\[
t_f=\sqrt{\frac{2H}{g}}.
\]
超人只有在 $T<t_f$ 时还有追赶机会。令他出发后的时间为
$\tau=t-T$，则
\[
y_s(\tau)=v_0\tau+\frac12g\tau^2.
\]
要使所需初速度最小，应把允许追赶的时间用到最长，即在临界情况下恰于女子落地时追上。此时
\[
\Delta=t_f-T
\]
是超人的飞行时间，并满足
\[
H=v_0\Delta+\frac12g\Delta^2.
\]
解得
\[
\boxed{v_{0,\min}=\frac{H-\frac12g(t_f-T)^2}{t_f-T}
=\frac{gT(2t_f-T)}{2(t_f-T)}}.
\]
第二个形式可由 $H=gt_f^2/2$ 化简得到。

用相对运动复核更能看清条件。超人出发时，女子已经领先
\[
s_0=\frac12gT^2
\]
并具有向下速度 $gT$。两人以后加速度相同，所以相对加速度为零；超人相对女子的接近速度恒为 $v_0-gT$。临界条件是
\[
(v_0-gT)(t_f-T)=\frac12gT^2,
\]
解得同一结果。题设允许临界状态，所以等号给出可取得的最小值；若另要求严格在落地前追上，则应取 $v_0>v_{0,\min}$，此时上式只是下确界。当 $T\to0$ 时该下界趋于零，当 $T\to t_f^-$ 时所需初速度趋于无穷，均符合直觉。
\end{xesolution}

\begin{origexamplebox}{恢复系数与无限次反弹}\label{ch02:ex:k02-2-19}
质点从高度 $H$ 自由下落，每次碰地后的反弹速率为碰前速率的 $e$ 倍，$0<e<1$。求从第一次反弹后算起到最终停止的总时间和总路程。
\par 题图见图~\ref{ch02:ref:question-figure-2-19}。
\KQuestionAT
\end{origexamplebox}
\Needspace{6\baselineskip}\begin{xesolution}
从高度 $H$ 自由下落到第一次碰地前，由
\[
u^2=2gH
\]
得碰前速率
\[
u=\sqrt{2gH}.
\]
恢复系数为 $e$ 表示每次碰地后的反弹速率等于碰前速率的 $e$ 倍。由于每次飞行都回到同一地面高度，下一次碰地前的速率又等于该次反弹初速，所以历次反弹初速度为
\[
eu,\ e^2u,\ e^3u,\ldots.
\]

若某次反弹初速度为 $w$，上升到最高点用时 $w/g$，再下落回地面也用时 $w/g$，故一轮飞行时间为
\[
\frac{2w}{g}.
\]
最高高度为 $w^2/(2g)$，上、下路程之和为
\[
2\frac{w^2}{2g}=\frac{w^2}{g}.
\]
题目从第一次反弹后开始计时和计路程，因此从 $w=eu$ 对应的项开始求和：
\BounceDiagram
\[
\begin{aligned}
T&=\frac{2u}{g}\sum_{n=1}^{\infty}e^n
=\frac{2u}{g}\frac{e}{1-e}\\
&=\boxed{\frac{2e\sqrt{2gH}}{g(1-e)}},
\end{aligned}
\]
\[
\begin{aligned}
S&=\frac{u^2}{g}\sum_{n=1}^{\infty}e^{2n}
=\frac{u^2}{g}\frac{e^2}{1-e^2}\\
&=\boxed{\frac{2He^2}{1-e^2}}.
\end{aligned}
\]
时间与速度成正比，所以按 $e$ 求和；路程与速度平方成正比，所以按 $e^2$ 求和，二者不能混用。检查极限：$e\to0$ 时第一次碰撞后不再反弹，$T,S\to0$；$e\to1^-$ 时能量几乎不损失，总时间和总路程都发散。若题目要求从最初释放算起，还应分别再加第一次下落时间 $\sqrt{2H/g}$ 和路程 $H$。
\end{xesolution}

\begin{origexamplebox}{逐次增强的回向加速度}\label{ch02:ex:k02-2-20}
质点以速度大小 $v_0$ 撞向一固定板。第一次及以后每次碰板时，速度均立即反向且大小不变。第一次反弹后，它以大小为 $a$、方向始终指向板的加速度运动；以后每完成一次碰板，回向加速度的大小变为前一次的 $k>1$ 倍。求第一次反弹后的总时间和总路程。
\end{origexamplebox}
\Needspace{6\baselineskip}\begin{xesolution}
把板的位置取为 $x=0$，离开板的方向取为正。第 $n$ 个“离板—回板”循环中，加速度大小为
\[
a_n=ak^{n-1}\qquad(n=1,2,\ldots),
\]
方向始终指向板，所以在 $x>0$ 区域实际加速度为 $-a_n$。每次碰板后，质点都以速度 $+v_0$ 离板。

在第 $n$ 个循环内，离板后经过时间 $\tau$ 的位置和速度为
\[
x=v_0\tau-\frac12a_n\tau^2,
\qquad
v=v_0-a_n\tau.
\]
令 $v=0$，得转向时刻 $\tau_1=v_0/a_n$；此时最远距离为
\[
x_{\max,n}=\frac{v_0^2}{2a_n}.
\]
令 $x=0$ 并舍去 $\tau=0$，得回到板的时刻
\[
\tau_n=\frac{2v_0}{a_n}.
\]
往返路程是最远距离的两倍：
\[
s_n=2x_{\max,n}=\frac{v_0^2}{a_n}.
\]
所以总时间和总路程分别为几何级数
\[
\begin{aligned}
T&=\frac{2v_0}{a}\sum_{n=0}^{\infty}k^{-n}
=\boxed{\frac{2v_0}{a}\frac{k}{k-1}},
\end{aligned}
\]
\[
\begin{aligned}
S&=\frac{v_0^2}{a}\sum_{n=0}^{\infty}k^{-n}
=\boxed{\frac{v_0^2}{a}\frac{k}{k-1}}.
\end{aligned}
\]
级数收敛依赖题设 $k>1$。单位检查：$v_0/a$ 是时间，$v_0^2/a$ 是长度。再检查极限：$k\to\infty$ 时，后续循环的加速度极大、耗时和路程趋零，结果只剩第一次循环的 $2v_0/a$ 与 $v_0^2/a$；$k\to1^+$ 时各循环几乎一样长，总和发散。
\end{xesolution}

\begin{origexamplebox}{周期运动的频闪取样}\label{ch02:ex:k02-2-21}
竖直真空玻璃管内，小球从距管底高度 $H$ 处由静止下落，与管底完全弹性碰撞后不停上下运动。照相机每隔固定时间 $T$ 拍摄一次，曝光时间可忽略；所有照片中小球均处于同一位置。求可能的 $H$，以及相应照片中小球距管底的可能高度。
\end{origexamplebox}
\Needspace{6\baselineskip}\begin{xesolution}
设从最高点落到底部的时间
\[
\tau=\sqrt{\frac{2H}{g}},
\]
完整运动周期为 $P=2\tau$。同一高度除转折点外在一个周期内出现两次，且两次相位关于最高点对称。

\begin{center}
\begin{tikzpicture}[x=1.25cm,y=2.05cm,>=Latex]
  \draw[axis,->] (-.15,0)--(4.35,0) node[right] {$t$};
  \draw[axis,->] (0,-.06)--(0,1.18) node[above] {$y/H$};
  \draw[curve,domain=0:2,samples=80] plot (\x,{1-(\x/2)^2});
  \draw[curve,domain=2:4,samples=80] plot (\x,{1-((4-\x)/2)^2});
  \draw[guide] (0,.75)--(4,.75);
  \fill (1,.75) circle (1.7pt);
  \fill (3,.75) circle (1.7pt);
  \node[below left=3pt] at (1,.75) {$t_1$};
  \node[below right=3pt] at (3,.75) {$t_2$};
  \draw[guide] (1,.75)--(1,.92) (3,.75)--(3,.92);
  \draw[<->,DeepBlue] (1,.92)--(3,.92) node[midway,above=1pt] {$\tau$};
  \foreach \x/\lab in {2/$\tau$,4/$P$}{\draw (\x,.025)--(\x,-.025) node[below] {\lab};}
  \node[left] at (0,1) {$1$};
\end{tikzpicture}
\figcap{一个周期内的高度—时间图。除最高点和碰撞点外，同一高度对应下降、上升两个相位；二者只有在图示特殊位置才相差半周期。}
\end{center}

为什么不会遗漏其他取样周期？把拍摄时刻按周期 $P$ 取模。若 $T/P$ 为无理数，取样相位会在整个周期内稠密分布，
高度显然不可能恒定；若 $T/P=p/q$ 为最简分数，一个循环会访问 $q$ 个不同相位。
而图中任一高度在一个周期内至多对应两个相位，所以只能有 $q=1$ 或 $q=2$。
它们分别对应相位完全不变和每次跨越半个周期，恰好就是下面两类。

若相机开始拍摄的相位未指定，能使无限序列始终停留在同一高度的取样只有两类：
\begin{enumerate}
  \item $T=nP=2n\tau$，$n=1,2,\ldots$。此时每次相位完全相同，
  \[
    \boxed{H=\frac{gT^2}{8n^2}},
  \]
  拍摄高度可为 $0\le y\le H$，由第一张照片的相位决定。
  \item $T=(2n+1)\tau$，$n=0,1,2,\ldots$。相位每次相差半个周期；要使两相位高度相同，它们必须位于上、下行关于中点对称的位置，即距最高点的等效时间均为 $\tau/2$。于是
  \[
    \boxed{H=\frac{gT^2}{2(2n+1)^2}},\qquad
    \boxed{y=H-\frac12g\left(\frac{\tau}{2}\right)^2=\frac34H}.
  \]
\end{enumerate}

第一类中，$T=nP$ 表示相邻两张照片之间恰好经过整数个完整周期，所以不仅高度相同，速度的大小和方向也完全相同。由
$T=2n\tau$ 得 $\tau=T/(2n)$，代入 $H=g\tau^2/2$，立即得到
$H=gT^2/(8n^2)$。由于首次曝光可以发生在周期内任意相位，照片中的高度也可以取闭区间 $[0,H]$ 内的任意值；一旦第一张照片确定，后续照片就重复同一状态。

第二类中，相邻照片的相位相差半个周期，速度方向相反，因而只能利用上升段和下降段在同一高度处的对称性。把两次经过该高度的时刻记为 $t_1,t_2$，图中有
$t_2-t_1=\tau=P/2$。两时刻分别位于下降、上升过程的中点，因此从最高点计的飞行时间都是 $\tau/2$，高度为
\[
 y=H-\frac12g\left(\frac{\tau}{2}\right)^2
  =H-\frac14\left(\frac12g\tau^2\right)=\frac34H.
\]
这里使用了 $H=g\tau^2/2$。两类结果的量纲均正确：$gT^2$ 是长度；当 $n$ 增大时，在固定拍摄间隔 $T$ 下允许的 $H$ 按 $1/n^2$ 缩小，与单程时间 $\tau\propto\sqrt H$ 相符。以上讨论还默认碰撞瞬间完成且完全弹性；若恢复系数小于 $1$，运动不再严格周期，结论不能直接沿用。

若题意另规定第一张照片就在释放瞬间拍摄，则只有第一类，且所有照片都拍到最高点 $y=H$。把“首次拍摄相位”说清楚，是这道题条件完整性的关键。
\end{xesolution}
}

\endgroup
